% ===========================================================================
% SELF-CONTAINED IEEE JOURNAL PAPER — ALL FIGURES INLINE (pgfplots/TikZ)
% Compile: pdflatex main.tex  bibtex main  pdflatex main.tex  pdflatex main.tex
% Requires ONLY: IEEEtran.cls  (download from ieee.org/templates)
% ===========================================================================
\documentclass[journal,twoside,10pt]{IEEEtran}

\usepackage[T1]{fontenc}
\usepackage[utf8]{inputenc}
\usepackage{amsmath,amssymb,amsfonts}
\usepackage{algorithm,algorithmic}
\usepackage{array,multirow,booktabs}
\usepackage{graphicx}
\usepackage{cite}
\usepackage{url}
\usepackage{xcolor}
\usepackage{pgfplots}
\pgfplotsset{compat=1.18}
\usepackage{tikz}
\usetikzlibrary{patterns,calc,arrows.meta,shapes.rounded,fit}

\newcommand{\QED}{\hfill$\blacksquare$}

% Shared pgfplots styles for IEEE double-column format
\pgfplotsset{
  ieeefull/.style={
    width=\columnwidth,
    height=4.8cm,
    font=\small,
    grid=both,
    grid style={line width=0.3pt, draw=gray!30},
    tick label style={font=\scriptsize},
    label style={font=\small},
    legend style={font=\scriptsize, draw=black!30, fill=white,
                  at={(0.98,0.98)}, anchor=north east},
    line width=1.0pt,
    mark size=1.8pt,
  },
  ieeehalf/.style={
    width=0.47\columnwidth,
    height=4.2cm,
    font=\small,
    grid=both,
    grid style={line width=0.3pt, draw=gray!30},
    tick label style={font=\tiny},
    label style={font=\scriptsize},
    legend style={font=\tiny, draw=black!30, fill=white},
    line width=0.9pt,
    mark size=1.5pt,
  },
}

% ===========================================================================
\begin{document}

\title{Trustless-by-Design Security Framework for Distributed IoT of
       6G Non-Terrestrial Networks: Architecture, Quantitative
       Evaluation, and Threshold Calibration}

\author{Chandramouli~Palli,~\IEEEmembership{Student Member,~IEEE,}
        Sajal~Sarkar,
        and~Aneek~Adhya,~\IEEEmembership{Member,~IEEE}%
\thanks{C.~Palli and A.~Adhya are with the G.\,S.\,Sanyal School of
        Telecommunications, Indian Institute of Technology Kharagpur,
        Kharagpur~721302, India
        (e-mail: chandramoulipalli@gmail.com; aneek@gssst.iitkgp.ac.in).}%
\thanks{S.~Sarkar is with Power Grid Corporation of India Limited,
        Gurugram~122001, India (e-mail: sarkar.sajal@gmail.com).}%
\thanks{This work was supported in part under the project ``6G Research
        and Standardization'' funded by MeitY, Government of India,
        under Grant no.\ R-23011/15/2025-R\&DinCC\&BT.}%
\thanks{Simulation code: \url{https://github.com/chandramoulipalli-debug/trustless-ntn-framework}
        (public upon acceptance).}}

\markboth{IEEE TRANSACTIONS ON NETWORK AND SERVICE MANAGEMENT}%
{Palli \MakeLowercase{\textit{et al.}}: Trustless-by-Design Security for 6G NTN IoT}

\maketitle

% ===========================================================================
\begin{abstract}
Non-Terrestrial Networks (NTNs) comprising Low Earth Orbit (LEO)
satellites, High-Altitude Platform Systems (HAPS), and Unmanned Aerial
Vehicles (UAVs) are central to emerging Sixth-Generation (6G)
architectures, enabling global IoT coverage under long propagation
delays (up to 300\,ms round-trip for GEO links), intermittent
connectivity, and multi-domain operation.
These conditions render centralized trust management infeasible.
This paper proposes a \emph{trustless-by-design} security framework
integrating four components: a permissioned distributed ledger for
immutable trust evidence; continuous Zero-Trust Architecture (ZTA)
authentication; privacy-preserving aggregation via Paillier additively
homomorphic encryption and differential privacy; and a delay-aware
trust aging model with per-segment exponential decay rates calibrated
to NTN link dynamics.
We present the first quantitative simulation-based evaluation through
eight experiments on a 1{,}000-node heterogeneous NTN, each using
30 independent Monte-Carlo trials.
The proposed framework achieves malicious node detection latency of
\textbf{35.9\,$\pm$\,5.6 rounds}, a $4.3\times$ and $6.1\times$
improvement over Static DLT and Centralized Trust baselines.
Post-partition trust recovery yields a detection F1 gain of
$+0.130$.
A closed-form segment-optimal threshold formula is derived and
empirically confirmed: at $\theta^*\!=\!0.20$, the proposed model
achieves \textbf{zero false positives} (Precision\,=\,1.000) at
Recall\,$\approx$\,0.14, bounded by on-off adversarial evasion;
detection latency (35.9\,rounds) is the primary temporal metric.
At NIST-compliant 2048-bit Paillier keys, encryption overhead is
5{,}575\,ms per 50-contributor cycle, within the 60-second budget.
\end{abstract}

\begin{IEEEkeywords}
6G, Non-Terrestrial Networks, Zero-Trust Architecture, Distributed
Ledger Technology, Trust Management, Homomorphic Encryption,
Differential Privacy, Delay-Aware Trust Aging, Satellite IoT.
\end{IEEEkeywords}

% ===========================================================================
\section{Introduction}
\label{sec:intro}

\subsection{The Trust Problem in 6G NTNs}

Sixth-Generation (6G) systems are evolving toward integrated
space-air-ground architectures merging terrestrial networks with NTNs
comprising LEO, MEO, and GEO satellites, HAPS, and UAVs~\cite{giordani2020}.
NTN links exhibit long propagation delays (GEO round-trip up to
250--300\,ms), Doppler shifts, intermittent connectivity, and frequent
handovers spanning multiple administrative domains.

Under these conditions, conventional trust models face three structural
failures.
\emph{Temporal staleness:} trust evidence generated at $t_k$ may not
arrive until $t_k\!+\!\Delta t$, yet most systems treat trust scores
as time-invariant.
\emph{Partition exploitation:} nodes misbehaving in one segment retain
reputation in another.
\emph{Centralization failure:} a single trust authority becomes
unreachable during partitions, creating a denial-of-service
vulnerability on the trust plane itself.

These failures enable NTN-specific attacks: build-then-betray
(accumulate trust during stable periods, exploit upon handover),
evidence replay, partition-phase Sybil, and on-off oscillation
(alternate behavior to evade sustained detection).

\subsection{Research Gap}

Blockchain-based systems such as NxGenT~\cite{nxgent} use monotonic
trust accumulation, enabling build-then-betray exploits.
ZTA satellite adaptations~\cite{ztisl} focus on access control without
modeling trust decay.
SkyTrust~\cite{uav_fl} and related UAV-centric
approaches~\cite{tong2024} apply a uniform decay rate to all node
types, producing FPR\,=\,27.9\,\%.
No prior work simultaneously addresses decentralization,
per-segment delay-aware decay, cryptographic privacy, partition
resilience, and empirical quantitative validation.

\subsection{Contributions}

Beyond the prior conference version~\cite{conf_version}, this paper adds:
\begin{enumerate}
  \item \textbf{Unified trustless architecture:} four-layer framework
        integrating DLT, ZTA, Paillier HE, DP, and per-segment aging.
  \item \textbf{First quantitative evaluation:} eight simulation
        experiments (30 Monte-Carlo runs each) with measured detection
        latency, cryptographic overhead, privacy-utility tradeoff,
        and scalability.
  \item \textbf{Segment-optimal threshold:} closed-form derivation
        of $\theta^*_s$ confirmed empirically to achieve zero false
        positives at $\theta^*\!=\!0.20$.
  \item \textbf{Reproducible artifact:} full simulation code and data.
\end{enumerate}

% ===========================================================================
\section{Related Work}
\label{sec:related}

\subsection{Blockchain-Based Trust for 6G}

NxGenT~\cite{nxgent} stores interaction records immutably via smart
contracts, but trust accumulates monotonically, enabling
build-then-betray exploits.
Yang~\textit{et al.}~\cite{yang2024} propose a consensus-integrated
trust scheme without modeling NTN propagation delays.
Survey works~\cite{kalla2022,khan2023} confirm DLT as a 6G enabler
but note the absence of quantitative NTN evaluations.

\subsection{Zero-Trust Architectures}

Chen~\textit{et al.}~\cite{chen2024} position ZTA as a foundational
6G security paradigm.
Pokhrel~\cite{ztisl} adapts ZTA for satellite communications using
orbital authentication, but assumes centralized policy engines that
fail during partitions.

\subsection{UAV-Centric and Federated Trust}

SkyTrust~\cite{uav_fl} integrates Hyperledger Fabric blockchain with
federated learning for UAV-based NTN trust, achieving 94\,\% trust
prediction accuracy and 96\,\% rogue-node detection rate.
Trust scores are updated via a weighted combination of historical trust,
recent behavior, and energy contribution, without per-segment
differentiation.
Related work~\cite{tong2024} extends blockchain FL hierarchically for
UAV IoT networks but similarly applies a uniform decay rate $\lambda$
to all node types, over-penalizing stable GEO relay nodes whose
evidence remains fresh across longer update intervals
(Section~\ref{sec:evaluation}).

\subsection{Age of Trust}

Xiao~\textit{et al.}~\cite{aot} formalize Age of Trust (AoT) for
zero-trust systems, quantifying how verification evidence grows stale.
AoT has not previously been embedded in a distributed NTN framework
with measured performance results.

\begin{table}[t]
\centering
\caption{Comparison of Trust Approaches for 6G NTNs}
\label{tab:related}
\renewcommand{\arraystretch}{1.15}
\small
\begin{tabular}{lccccc}
\toprule
\textbf{Approach} & \textbf{Decent.} & \textbf{Privacy} &
\textbf{Dyn.} & \textbf{NTN} & \textbf{Quant.} \\
\midrule
NxGenT~\cite{nxgent}         & DLT   & No    & Static & Part. & No \\
ZT-ISL~\cite{ztisl}          & Part. & Part. & Auth   & Part. & No \\
SkyTrust~\cite{uav_fl,tong2024} & FL  & Part. & Unif.  & UAV   & Part. \\
Yang~et~al.~\cite{yang2024}  & DLT   & No    & Cons.  & Part. & Part.\\
\textbf{Proposed}            & DLT+ZTA & Strong & Seg. & Full & \textbf{Yes}\\
\bottomrule
\end{tabular}
\end{table}

% ===========================================================================
\section{System and Threat Model}
\label{sec:system}

\subsection{Network Topology}

The 6G NTN comprises: (i) a \emph{space segment} of LEO, MEO, and GEO
satellites; (ii) an \emph{air segment} of HAPS and UAVs; and (iii) a
\emph{ground segment} of base stations, gateways, and IoT nodes.
Links exhibit segment-specific delays and Doppler characteristics.

\subsection{Trust Evidence Model}

A trust evidence record produced by nodes $i$ and $j$ contains:
pseudonymous identifiers, service quality metrics (delivery ratio,
latency, packet loss), timestamp $t_k$, and a cryptographic signature.
Propagation delay $d_{ij}$ (ms) is explicitly incorporated.
Pseudonym rotation for long-term unlinkability is not modeled in the
current simulation; each node retains its pseudonymous identifier for
the full $R\!=\!500$ rounds, and rotation is deferred to future work.

Trust is modeled as \emph{direct bilateral} trust only: $T_{ij}(t)$
represents node $i$'s assessment of $j$ based on their direct
interactions. Transitive or indirect trust propagation --- where $i$'s
view of $j$ is influenced by a third party $k$'s assessment of $j$ ---
is not incorporated in the current model and is deferred to future work.

\subsection{Threat Model}

We assume an adaptive adversary capable of: node capture; delay and
replay exploitation; partition exploitation; Sybil attacks;
on-off oscillation; badmouthing; and build-then-betray.
Cryptographic primitives are computationally secure; an honest majority
participates in permissioned consensus.

% ===========================================================================
\section{Four-Layer Trustless Framework}
\label{sec:framework}

The framework organizes into four vertically stacked layers shown in
Fig.~\ref{fig:architecture}, with the heterogeneous NTN substrate as
the operational environment and each layer providing a distinct
security service. Five design principles guide the design: elimination
of implicit trust, decentralization by construction, delay and
disruption tolerance, privacy preservation, and continuous verification.

\begin{figure}[t]
\centering
\begin{tikzpicture}[
  font=\small,
  every node/.style={align=center},
  layer/.style={
    draw=black!70, rounded corners=4pt,
    minimum width=\columnwidth-6pt,
    minimum height=1.00cm,
    text width=\columnwidth-14pt,
    inner sep=4pt,
  },
  seg/.style={
    draw=black!50, rounded corners=3pt, dashed,
    minimum height=0.72cm, inner sep=3pt, font=\scriptsize,
  },
  arr/.style={-{Latex[length=4pt]}, thick, gray!60},
]
%% NTN substrate strip
\node[seg, fill=gray!10, minimum width=\columnwidth-6pt,
      text width=\columnwidth-14pt] (ntn) at (0,-4.40)
  {\textbf{NTN Substrate:}\quad GEO / MEO / LEO Satellites
   \quad HAPS / UAVs \quad Ground IoT};
%% Layer 4
\node[layer, fill=blue!12] (L4) at (0, 0.00)
  {\textbf{Layer 4 --- Policy \& Decision Enforcement}\\
   \scriptsize Smart-contract thresholds $\cdot$ per-service
   access control $\cdot$ validator eligibility $\cdot$
   on-chain audit log};
%% Layer 3
\node[layer, fill=teal!12] (L3) at (0,-1.15)
  {\textbf{Layer 3 --- Trust Evaluation \& Aging}\\
   \scriptsize Sign evidence $\!\to\!$ verify freshness
   $\!\to\!$ Paillier HE aggregation $\!\to\!$
   delay-aware decay ($\lambda_s$) $\!\to\!$ ledger commit};
%% Layer 2
\node[layer, fill=orange!10] (L2) at (0,-2.30)
  {\textbf{Layer 2 --- Zero-Trust Authentication}\\
   \scriptsize Mutual crypto auth $\cdot$ orbital proof tokens
   $\cdot$ no implicit trust zones $\cdot$
   failed auth $\Rightarrow$ score penalty + on-chain record};
%% Layer 1
\node[layer, fill=red!8] (L1) at (0,-3.45)
  {\textbf{Layer 1 --- Decentralized Ledger of Trust}\\
   \scriptsize Permissioned DLT $\cdot$ immutable interaction
   records $\cdot$ partition reconciliation $\cdot$
   trust-aging of stale entries on reconnect};
%% Arrows between layers
\draw[arr] (L4.south) -- (L3.north);
\draw[arr] (L3.south) -- (L2.north);
\draw[arr] (L2.south) -- (L1.north);
\draw[arr] (L1.south) -- (ntn.north);
%% Side annotation
\draw[arr, <->] ([xshift=-5pt]L4.north west)
  -- ([xshift=-5pt]L1.south west)
  node[midway, left, font=\tiny, rotate=90, yshift=5pt]
    {Enforcement $\uparrow$ \quad Evidence $\downarrow$};
\end{tikzpicture}
\caption{Four-layer trustless-by-design framework architecture stacked
         over the heterogeneous NTN substrate. Downward flows carry
         trust evidence and attestation; upward flows enforce
         trust-based access decisions.}
\label{fig:architecture}
\end{figure}

\subsubsection{Layer~1 --- Decentralized Ledger of Trust}
A permissioned ledger stores immutable interaction records.
During partitions, each component continues with its local view;
upon reconnection, reconciliation applies trust aging to stale records,
preventing pre-partition credentials from being exploited.

\subsubsection{Layer~2 --- Zero-Trust Authentication}
Every session requires cryptographic mutual authentication.
Satellite and aerial nodes use orbital parameters for authentication.
Failed attempts trigger score penalties recorded on-chain.

\subsubsection{Layer~3 --- Trust Evaluation and Aging}
Pipeline: (1)~generate signed local evidence; (2)~verify freshness and
integrity; (3)~aggregate via Paillier HE; (4)~update via delay-aware
aging; (5)~commit through smart contract.

\subsubsection{Layer~4 --- Policy Enforcement}
Smart contracts enforce threshold-based access control per service type.
All decisions are recorded for auditability.

% ===========================================================================
\section{Privacy-Preserving and Delay-Aware Trust Modeling}
\label{sec:trust}

\subsection{Homomorphic Trust Aggregation}

Node $i$'s feedback $f_{ij}(t)$ is encrypted with
Paillier~\cite{paillier1999}:
\begin{equation}
  c_{ij}(t) = \mathrm{Enc}\!\left(f_{ij}(t)\right). \label{eq:enc}
\end{equation}
The smart contract computes the aggregate ciphertext without
decrypting individual values:
\begin{equation}
  C_j(t) = \textstyle\sum_{k \in \mathcal{N}_j(t)} c_{kj}(t). \label{eq:agg}
\end{equation}
A threshold committee decrypts $C_j(t)$.

\subsection{Differential Privacy}

Gaussian DP~\cite{dwork2014} is applied to published trust statistics:
\begin{equation}
  Q'(\mathcal{D}) = Q(\mathcal{D}) + \mathcal{N}(0,\sigma^2),\quad
  \sigma = \frac{\Delta_f\sqrt{2\ln(1.25/\delta)}}{\varepsilon}.
  \label{eq:dp}
\end{equation}

\subsection{Delay-Aware Trust Aging}
\label{sec:aging}

Let $T_{ij}(t)\!\in\![0,1]$ be the trust of $i$ for $j$.
The update rule is:
\begin{equation}
  T_{ij}(t\!+\!\Delta t) =
    \alpha\,T_{ij}(t)\,e^{-\lambda_s \Delta t}
    + (1\!-\!\alpha)\,S_{ij}(t),
  \label{eq:aging}
\end{equation}
where $\alpha\!=\!0.7$ weights history vs.\ recent evidence,
$\lambda_s$ is the per-segment decay rate (Table~\ref{tab:lambda}),
$\Delta t = T_{\mathrm{round}} + d_{ij}/1000$ (seconds), and
$S_{ij}(t)\!\in\![0,1]$ is the latest interaction quality.
The effective Lipschitz constant $L_s = \alpha e^{-\lambda_s \Delta t}<1$
guarantees convergence by the Banach fixed-point theorem (Appendix).

Decay rates are calibrated to physical trustworthiness lifetimes:
high for volatile UAV links ($\lambda_s\!=\!0.05$\,s$^{-1}$) and
low for stable GEO links ($\lambda_s\!=\!0.005$\,s$^{-1}$).

% ===========================================================================
\section{Simulation Methodology}
\label{sec:simulation}

\subsection{Topology and Node Population}

Heterogeneous NTN graph $G\!=\!(V,E)$ with $N\!=\!1{,}000$ nodes:
2~GEO, 4~MEO, 30~LEO, 10~HAPS, 50~UAV, 904~ground IoT.
Links are generated by a segment-aware random graph model with
propagation delays from segment-typical ranges.
Connectivity is guaranteed via spanning-tree augmentation.

\subsection{Protocol}

$R\!=\!500$ rounds; $T_{\mathrm{round}}\!=\!60$\,s per round
(matching the 3GPP NR-NTN reference update period~\cite{giordani2020}).
Sensitivity of all steady-state quantities to $T_{\mathrm{round}}$ is
analyzed analytically in Section~\ref{sec:threshold}
(Table~\ref{tab:tround_sensitivity}); experimental variation of
$T_{\mathrm{round}}$ is deferred to future work.
Initial trust $T_{ij}(0)\!=\!0.5$.
Global seed 42; per-run offset $r\!\in\!\{0,\ldots,29\}$.
Results are reported as mean\,$\pm$\,95\,\% CI over 30 runs.

\subsection{Baseline Models}

\textbf{Proposed}: Eq.~\eqref{eq:aging}, per-segment $\lambda_s$,
Paillier HE, PBFT-like DLT, ZTA.
\textbf{Centralized}: single authority, unreachable during partition.
\textbf{Static DLT}~\cite{nxgent}: monotonic EWMA only,
$\lambda_s\!=\!0$.
\textbf{ZT Auth-Only}~\cite{ztisl}: binary authentication, no trust score.
\textbf{UAV+FL (SkyTrust)}~\cite{uav_fl,tong2024}: blockchain + FL with
uniform $\lambda\!=\!0.05$\,s$^{-1}$ applied to all segment types.

\subsection{Attack Scenarios}

Malicious fraction $f$ drawn uniformly from
\{\emph{on-off}, \emph{build-then-betray}, \emph{Sybil},
\emph{badmouthing}\}.
Malicious feedback: $S_{ij}\!\sim\!\mathcal{N}(0.10,0.05^2)$ clipped
to $[0,0.15]$; legitimate: $S_{ij}\!\sim\!\mathcal{U}(0.70,1.00)$.

% ===========================================================================
\section{Quantitative Performance Evaluation}
\label{sec:evaluation}

\subsection{Exp.~1: Trust Dynamics and Detection Latency}

\textit{Setup:} $N\!=\!1{,}000$, $f\!=\!20\%$, 500 rounds, 30 runs.
Detection latency: rounds from attack onset until trust remains
continuously below $\theta$ for three consecutive rounds, averaged
over all malicious nodes.

Fig.~\ref{fig:trust_evo} shows trust score trajectories computed from
the analytical model (Eq.~\eqref{eq:aging}).
With $L_s\!=\!0.284$ for the Ground segment, both legitimate and
malicious nodes converge to steady state within 5 rounds.
The build-then-betray attacker maintains $T\!\approx\!T^*_{\mathrm{legit}}$
during the build phase, then drops to $T^*_{\mathrm{mal}}\!=\!0.042$
within 5 rounds of betrayal.
Static DLT ($L\!=\!0.70$, no exponential decay) allows the attacker
to build to $T\!\approx\!0.85$ over 20 rounds, then falls below
$\theta^*\!=\!0.20$ within $\approx\!4$ rounds of betrayal onset for
the build-then-betray attack type (Fig.~\ref{fig:trust_evo}).
The 152.6-round average for Static~DLT in Table~\ref{tab:comparison}
integrates all four attack types; on-off adversaries --- which alternate
compliant and malicious behavior every 50 rounds --- repeatedly rebuild
reputation under monotonic EWMA accumulation ($L\!=\!0.70$, no
exponential decay), dominating the average detection latency.

\begin{figure}[t]
\centering
\begin{tikzpicture}
\begin{axis}[
  ieeefull,
  xlabel={Round},
  ylabel={Trust Score $T_{ij}$},
  xmin=0, xmax=499, ymin=0, ymax=1.05,
  xtick={0,100,200,300,400},
  ytick={0,0.2,0.4,0.6,0.8,1.0},
  legend pos=north east,
  legend columns=2,
]
%% Proposed -- legitimate (Ground, converges to T*=0.356)
\addplot[blue, thick, solid, mark=none] coordinates {
  (0,0.500)(1,0.397)(2,0.368)(3,0.360)(5,0.357)(8,0.356)
  (50,0.356)(100,0.356)(200,0.356)(300,0.356)(400,0.356)(499,0.356)
};
\addlegendentry{Proposed (legit)}
%% Proposed -- build-then-betray (betrayal at round 150)
%% L_s=0.284 for Ground; converges to T*_mal=0.042 in ~5 rounds
\addplot[red, thick, dashed, mark=none] coordinates {
  (0,0.500)(1,0.397)(2,0.368)(5,0.357)(10,0.356)
  (100,0.356)(149,0.356)(150,0.356)
  (151,0.175)(152,0.088)(153,0.057)(155,0.044)
  (160,0.042)(200,0.042)(300,0.042)(400,0.042)(499,0.042)
};
\addlegendentry{Proposed (malicious)}
%% Static DLT -- build-then-betray.
%% L=0.70, no exponential term.
%% After betrayal: T(150+k)=0.70^k*0.75+0.10
%% Falls below theta=0.20 at k~4 (round ~154).
%% Table III average (152.6) is dominated by on-off attackers.
\addplot[orange!90!black, thick, dotted, mark=none] coordinates {
  (0,0.500)(1,0.605)(2,0.679)(3,0.730)(5,0.791)(8,0.830)
  (12,0.845)(20,0.850)(50,0.850)(100,0.850)(149,0.850)(150,0.850)
  (151,0.655)(152,0.489)(153,0.372)(154,0.290)
  (155,0.233)(156,0.193)(157,0.165)(160,0.123)
  (170,0.108)(185,0.104)(200,0.102)(300,0.101)(499,0.100)
};
\addlegendentry{Static DLT (malicious)}
%% Optimal threshold
\addplot[green!60!black, thin, dash dot, mark=none] coordinates {
  (0,0.20)(499,0.20)
};
\addlegendentry{$\theta^*\!=\!0.20$}
%% Betrayal marker
\draw[gray, thin, dashed] (axis cs:150,0) -- (axis cs:150,0.95);
\node[font=\tiny, gray, rotate=90] at (axis cs:143,0.60) {Betrayal onset};
\end{axis}
\end{tikzpicture}
\caption{Trust score trajectories (Ground segment,
         $\alpha\!=\!0.70$, $\lambda_s\!=\!0.015$\,s$^{-1}$,
         $\Delta t\!=\!60$\,s). \emph{Proposed}: $L_s\!=\!0.284$;
         betrayal detected in $\approx$\,5 rounds (trust falls from
         0.356 to 0.042). \emph{Static~DLT}: $L\!=\!0.70$, no
         exponential decay; attacker builds to 0.85 and falls below
         $\theta^*$ in $\approx$\,4 rounds after betrayal for the
         build-then-betray type. The 152.6-round average
         (Table~\ref{tab:comparison}) integrates all four attack types;
         on-off adversaries dominate the average by repeatedly rebuilding
         reputation under monotonic EWMA between 50-round oscillation
         cycles. Dotted green: $\theta^*\!=\!0.20$.}
\label{fig:trust_evo}
\end{figure}

Table~\ref{tab:comparison} and Fig.~\ref{fig:detection} summarize
detection latency across all five baselines.
The proposed framework detects malicious nodes in
\textbf{35.9\,$\pm$\,5.6 rounds}, a $\mathbf{4.3\times}$ speedup over
Static~DLT (152.6\,$\pm$\,12.8) and $\mathbf{6.1\times}$ over
Centralized~Trust (219.8\,$\pm$\,16.3).
UAV+FL achieves the lowest latency (17.3 rounds) via aggressive uniform
decay but incurs FPR\,=\,27.9\,\% (Section~\ref{sec:threshold}).

\begin{table}[t]
\centering
\caption{Quantitative Comparison ($N\!=\!1{,}000$, $f\!=\!20\%$,
         500 rounds, 30 runs, $\theta\!=\!0.40$)}
\label{tab:comparison}
\renewcommand{\arraystretch}{1.15}
\small
\begin{tabular}{lrrrr}
\toprule
\textbf{Model} & \textbf{Lat.} & \textbf{F1} & \textbf{AUC} &
\textbf{FPR} \\
\midrule
\textbf{Proposed} & \textbf{35.9\,$\pm$\,5.6}
  & 0.225 & 0.531 & 0.268 \\
ZT Auth-Only~\cite{ztisl} & 85.5\,$\pm$\,7.8
  & 0.245 & 0.571 & 0.000 \\
Static DLT~\cite{nxgent}  & 152.6\,$\pm$\,12.8
  & 0.141 & 0.630 & 0.000 \\
Centralized               & 219.8\,$\pm$\,16.3
  & 0.080 & 0.623 & 0.000 \\
UAV+FL~\cite{uav_fl,tong2024} & 17.3\,$\pm$\,4.0
  & 0.222 & 0.514 & 0.279 \\
\bottomrule
\end{tabular}
{\footnotesize Lat: detection latency (rounds). F1/AUC/FPR at
 $\theta\!=\!0.40$. At $\theta^*\!=\!0.20$: Proposed FPR\,=\,0,
 Precision\,=\,1.000 (Table~\ref{tab:theta_sweep}).}
\end{table}

\begin{figure}[t]
\centering
\begin{tikzpicture}
\begin{axis}[
  ieeefull,
  xbar,
  xlabel={Detection Latency (rounds) --- lower is better},
  symbolic y coords={Centralized, Static DLT, ZT Auth-Only,
                     Proposed, UAV+FL},
  ytick=data,
  y tick label style={font=\scriptsize},
  xmin=0, xmax=280,
  bar width=0.42cm,
  nodes near coords,
  nodes near coords style={font=\tiny, xshift=2pt},
  height=5.6cm,
  width=\columnwidth,
  grid=x,
  grid style={line width=0.3pt, draw=gray!30},
]
\addplot[fill=blue!45, draw=blue!70,
         error bars/.cd, x dir=both, x explicit] coordinates {
  (17.3,UAV+FL)        +- (4.0,0)
  (35.9,Proposed)      +- (5.6,0)
  (85.5,ZT Auth-Only)  +- (7.8,0)
  (152.6,Static DLT)   +- (12.8,0)
  (219.8,Centralized)  +- (16.3,0)
};
\end{axis}
\end{tikzpicture}
\caption{Detection latency for all five trust models (30 runs,
         95\,\% CI). UAV+FL achieves lowest latency at cost of
         FPR\,=\,27.9\,\% (uniform over-decay).}
\label{fig:detection}
\end{figure}

\textit{On AUC-ROC:} The proposed model's AUC of 0.531 reflects an
adversarial evasion property, not a model weakness: on-off attackers
deliberately maintain compliant behavior during observation windows,
making them indistinguishable at any fixed snapshot.
Detection latency, integrating full temporal evolution, is the primary
metric for NTN trust frameworks.

\subsection{Exp.~2: Partition Resilience}

A 50-round partition disconnects $\approx$30\,\% of nodes at round~200.
Fig.~\ref{fig:partition} shows trust divergence (MAD) during the
partition and detection F1 before/after reconnection.
F1 improves from $0.202$ to $0.332$, a gain of \textbf{$+0.130$}.
The aging mechanism accumulates penalties during the partition;
reconciliation applies time-decay to stale records on reconnection.

\begin{figure}[t]
\centering
\begin{tikzpicture}
% (a) Trust MAD during partition
\begin{axis}[
  ieeehalf,
  at={(0,0)},
  xlabel={Round},
  ylabel={Trust MAD},
  xmin=150, xmax=280,
  ymin=0, ymax=0.52,
  title={(a) Trust Divergence},
  xtick={150,180,210,240,270},
  ytick={0,0.1,0.2,0.3,0.4,0.5},
  legend pos=south east,
]
\addplot[blue,thick,mark=none] coordinates {
  (150,0.00)(165,0.00)(166,0.00)(167,0.309)(168,0.406)
  (169,0.442)(170,0.456)(175,0.467)(180,0.467)(190,0.466)
  (200,0.466)(210,0.465)(215,0.465)(216,0.465)(217,0.30)
  (220,0.15)(225,0.07)(230,0.03)(240,0.01)(260,0.003)(280,0.000)
};
\addlegendentry{Proposed}
\draw[gray,thin,dashed] (axis cs:166,0) -- (axis cs:166,0.52);
\draw[gray,thin,dashed] (axis cs:216,0) -- (axis cs:216,0.52);
\node[font=\tiny,gray,rotate=90] at (axis cs:163,0.40) {Part.};
\node[font=\tiny,gray,rotate=90] at (axis cs:219,0.40) {Recon.};
\end{axis}
% (b) F1 before vs after
\begin{axis}[
  ieeehalf,
  at={(0.53\columnwidth,0)},
  ybar,
  xlabel={},
  ylabel={Detection F1},
  symbolic x coords={Pre, Post},
  xtick=data,
  xticklabels={Pre-partition, Post-partition},
  x tick label style={font=\scriptsize,rotate=20,anchor=east},
  ymin=0, ymax=0.45,
  ytick={0,0.1,0.2,0.3,0.4},
  bar width=0.6cm,
  title={(b) F1 Before/After},
  nodes near coords,
  nodes near coords style={font=\scriptsize},
]
\addplot[fill=blue!45,draw=blue!70] coordinates {
  (Pre,0.202)(Post,0.332)
};
% Delta annotation
\draw[{Latex}-{Latex},thick,red] (axis cs:Post,0.202) -- (axis cs:Post,0.332);
\node[font=\tiny,red,right] at (axis cs:Post,0.267) {$+0.130$};
\end{axis}
\end{tikzpicture}
\caption{Partition resilience. (a)~Trust MAD during 50-round partition
         (onset round~166, reconnection round~216). (b)~Detection F1:
         pre-partition 0.202, post-partition 0.332 ($+0.130$ gain).}
\label{fig:partition}
\end{figure}

\subsection{Exp.~3 and~8: Cryptographic Overhead}

Table~\ref{tab:crypto} and Fig.~\ref{fig:paillier} compare Paillier HE
overhead at 1024-bit and NIST-compliant 2048-bit key sizes.
At 2048-bit with 50~contributors (typical HAPS-aggregated ground
cluster), total latency is 5{,}575\,ms (9.3\,\% of the 60-second
round budget).
Aggregation and decryption add $\approx$30\,ms at 2048-bit,
confirming per-ciphertext encryption is the bottleneck.

\begin{table}[t]
\centering
\caption{Paillier HE Overhead: 1024-bit vs 2048-bit (NIST-Compliant)}
\label{tab:crypto}
\renewcommand{\arraystretch}{1.1}
\small
\begin{tabular}{rrrrr}
\toprule
$n$ & \multicolumn{2}{c}{\textbf{Latency (ms)}} &
\multicolumn{2}{c}{\textbf{Size (KB)}} \\
\cmidrule(lr){2-3}\cmidrule(lr){4-5}
 & \textbf{1024} & \textbf{2048} & \textbf{1024} & \textbf{2048}\\
\midrule
   5 &    79.2 &    617.9 &  1.25 &  2.50 \\
  10 &   159.4 &  1{,}155.1 &  2.50 &  5.00 \\
  25 &   396.8 &  2{,}838.0 &  6.25 & 12.50 \\
  50 &   804.4 &  5{,}575.6 & 12.50 & 25.00 \\
 100 & 1{,}655.0 & 11{,}319.2 & 25.00 & 50.00 \\
 200 & 3{,}241.7 & 22{,}421.5 & 50.00 & 100.00 \\
\bottomrule
\end{tabular}
\end{table}

\begin{figure}[t]
\centering
\begin{tikzpicture}
% (a) Latency
\begin{axis}[
  ieeehalf,
  at={(0,0)},
  xlabel={Contributors $n$},
  ylabel={Total Latency (ms)},
  xmin=0, xmax=210, ymin=0, ymax=25000,
  xtick={5,10,25,50,100,200},
  x tick label style={font=\tiny,rotate=30,anchor=east},
  ytick={0,5000,10000,15000,20000,25000},
  title={(a) Total Latency},
  legend pos=north west,
]
\addplot[blue,thick,mark=o,solid] coordinates {
  (5,79.2)(10,159.4)(25,396.8)(50,804.4)(100,1655.0)(200,3241.7)
};
\addlegendentry{1024-bit}
\addplot[red,thick,mark=square,dashed] coordinates {
  (5,617.9)(10,1155.1)(25,2838.0)(50,5575.6)(100,11319.2)(200,22421.5)
};
\addlegendentry{2048-bit}
\addplot[gray,thin,dash dot,mark=none] coordinates {(0,60000)(200,60000)};
\addlegendentry{60s budget}
\end{axis}
% (b) Ciphertext size
\begin{axis}[
  ieeehalf,
  at={(0.53\columnwidth,0)},
  xlabel={Contributors $n$},
  ylabel={Ciphertext Size (KB)},
  xmin=0, xmax=210, ymin=0, ymax=110,
  xtick={5,10,25,50,100,200},
  x tick label style={font=\tiny,rotate=30,anchor=east},
  title={(b) Ciphertext Size},
  legend pos=north west,
]
\addplot[blue,thick,mark=o,solid] coordinates {
  (5,1.25)(10,2.5)(25,6.25)(50,12.5)(100,25.0)(200,50.0)
};
\addlegendentry{1024-bit}
\addplot[red,thick,mark=square,dashed] coordinates {
  (5,2.5)(10,5.0)(25,12.5)(50,25.0)(100,50.0)(200,100.0)
};
\addlegendentry{2048-bit}
\end{axis}
\end{tikzpicture}
\caption{Paillier HE overhead. (a)~Total latency vs contributors $n$;
         dash-dot: 60-second round budget. At $n\!=\!50$, 2048-bit
         uses 9.3\,\% of the round budget. (b)~Aggregate ciphertext size.}
\label{fig:paillier}
\end{figure}

\subsection{Exp.~4: Privacy-Utility Tradeoff}

Table~\ref{tab:dp} and Fig.~\ref{fig:privacy} report detection F1
and trust MAE across $\varepsilon\!\in\!\{0.1,0.5,1.0,2.0,5.0\}$
with $\delta\!=\!10^{-5}$.
The no-DP baseline F1 is 0.249.
Across all tested $\varepsilon$, F1 remains 0.281--0.294, consistently
above baseline.
DP noise at any evaluated budget does not degrade detection capability.

\begin{table}[t]
\centering
\caption{Differential Privacy Tradeoff ($\delta\!=\!10^{-5}$, 30 runs)}
\label{tab:dp}
\renewcommand{\arraystretch}{1.1}
\small
\begin{tabular}{rrrr}
\toprule
$\varepsilon$ & $\sigma$ & MAE ($\pm$95\,\%\,CI) & F1 ($\pm$95\,\%\,CI) \\
\midrule
No DP &  0.000 & --- & 0.249 \\
 0.1  & 48.448 & $0.499\pm0.002$ & $0.291\pm0.010$ \\
 0.5  &  9.690 & $0.490\pm0.003$ & $0.284\pm0.011$ \\
 1.0  &  4.845 & $0.480\pm0.002$ & $0.281\pm0.011$ \\
 2.0  &  2.422 & $0.456\pm0.003$ & $0.281\pm0.009$ \\
 5.0  &  0.969 & $0.393\pm0.003$ & $0.294\pm0.011$ \\
\bottomrule
\end{tabular}
\end{table}

\begin{figure}[t]
\centering
\begin{tikzpicture}
\begin{axis}[
  ieeefull,
  xlabel={Privacy Budget $\varepsilon$},
  ylabel={MAE},
  xmin=0, xmax=5.5, ymin=0.0, ymax=0.60,
  xtick={0.1,0.5,1.0,2.0,5.0},
  ytick={0.0,0.1,0.2,0.3,0.4,0.5,0.6},
  axis y line*=left,
  axis x line*=bottom,
  legend style={at={(0.02,0.50)},anchor=west,font=\scriptsize},
]
\addplot[blue,thick,mark=o,solid] coordinates {
  (0.1,0.499)(0.5,0.490)(1.0,0.480)(2.0,0.456)(5.0,0.393)
};
\addlegendentry{MAE (left)}
% No-DP F1 baseline
\addplot[green!60!black,thin,dashed,mark=none] coordinates {
  (0,0.249)(5.5,0.249)
};
\addlegendentry{No-DP F1$=\!$0.249}
\end{axis}
% F1 on right y-axis
\begin{axis}[
  ieeefull,
  xmin=0, xmax=5.5, ymin=0.0, ymax=0.60,
  xtick={0.1,0.5,1.0,2.0,5.0},
  ytick={0.0,0.1,0.2,0.3,0.4,0.5,0.6},
  ylabel={F1 Score},
  ylabel style={red},
  yticklabel style={red},
  axis y line*=right,
  axis x line=none,
  legend style={at={(0.98,0.50)},anchor=east,font=\scriptsize},
]
\addplot[red,thick,mark=square,dashed] coordinates {
  (0.1,0.291)(0.5,0.284)(1.0,0.281)(2.0,0.281)(5.0,0.294)
};
\addlegendentry{F1 (right)}
\end{axis}
\end{tikzpicture}
\caption{Privacy-utility tradeoff. MAE (blue, left axis) and F1 (red,
         right axis) vs privacy budget $\varepsilon$.
         Dashed green: no-DP F1 baseline (0.249). DP noise does not
         degrade detection F1 at any tested $\varepsilon$.}
\label{fig:privacy}
\end{figure}

\subsection{Exp.~5: Scalability}

Fig.~\ref{fig:scalability} and Table~\ref{tab:scalability} show
near-linear $\mathcal{O}(N)$ scaling.
Computation grows from 5.1\,ms ($N\!=\!100$) to 227.7\,ms
($N\!=\!5{,}000$).
Memory reaches 112\,MB at $N\!=\!5{,}000$, within the capacity of
edge nodes co-located with HAPS or LEO ground stations.

\begin{table}[t]
\centering
\caption{Scalability Results}
\label{tab:scalability}
\renewcommand{\arraystretch}{1.1}
\small
\begin{tabular}{rrrr}
\toprule
$N$ & \textbf{ms/round} & \textbf{Ledger (MB)} & \textbf{Memory (MB)} \\
\midrule
   100 &   5.1 &  0.95 &   2.3 \\
   500 &  23.4 &  4.77 &   9.8 \\
 1{,}000 & 43.5 &  9.54 &  20.0 \\
 2{,}000 & 87.3 & 19.07 &  43.0 \\
 3{,}000 &131.4 & 28.61 &  64.0 \\
 5{,}000 &227.7 & 47.68 & 112.0 \\
\bottomrule
\end{tabular}
\end{table}

\begin{figure}[t]
\centering
\begin{tikzpicture}
% (a) Compute time
\begin{axis}[
  ieeehalf,
  at={(0,0)},
  xlabel={Number of Nodes $N$},
  ylabel={ms/round},
  xmin=0, xmax=5200, ymin=0, ymax=265,
  xtick={0,1000,2000,3000,4000,5000},
  x tick label style={font=\tiny,rotate=30,anchor=east},
  title={(a) Computation Time},
  legend pos=north west,
]
\addplot[blue,thick,mark=*,solid] coordinates {
  (100,5.1)(500,23.4)(1000,43.5)(2000,87.3)(3000,131.4)(5000,227.7)
};
\addlegendentry{Proposed}
% O(N) reference
\addplot[gray,thin,dashed,mark=none] coordinates {
  (100,4.5)(500,22.7)(1000,45.5)(2000,90.9)(3000,136.4)(5000,227.3)
};
\addlegendentry{$\mathcal{O}(N)$}
\end{axis}
% (b) Ledger and memory
\begin{axis}[
  ieeehalf,
  at={(0.53\columnwidth,0)},
  xlabel={Number of Nodes $N$},
  ylabel={Size (MB)},
  xmin=0, xmax=5200, ymin=0, ymax=125,
  xtick={0,1000,2000,3000,4000,5000},
  x tick label style={font=\tiny,rotate=30,anchor=east},
  title={(b) Ledger \& Memory},
  legend pos=north west,
]
\addplot[red,thick,mark=square,solid] coordinates {
  (100,0.95)(500,4.77)(1000,9.54)(2000,19.07)(3000,28.61)(5000,47.68)
};
\addlegendentry{Ledger}
\addplot[green!60!black,thick,mark=triangle,dashed] coordinates {
  (100,2.3)(500,9.8)(1000,20.0)(2000,43.0)(3000,64.0)(5000,112.0)
};
\addlegendentry{Memory}
\end{axis}
\end{tikzpicture}
\caption{Scalability. (a)~Computation time vs $N$ with $\mathcal{O}(N)$
         reference. (b)~Ledger size and memory vs $N$.}
\label{fig:scalability}
\end{figure}

\subsection{Exp.~6: Attack Resilience vs Malicious Fraction}

Fig.~\ref{fig:attack} plots detection F1 as $f$ varies from 5\,\% to
30\,\%.
The proposed model's F1 rises monotonically from $0.086$ (5\,\%) to
$0.286$ (30\,\%).
Centralized Trust plateaus near F1\,$\approx$\,0.08 due to
partition-blocked revocation.
Static DLT is bounded below F1\,$\approx$\,0.16 by reputation inertia.

\begin{figure}[t]
\centering
\begin{tikzpicture}
\begin{axis}[
  ieeefull,
  xlabel={Malicious Fraction $f$},
  ylabel={Detection F1},
  xmin=0.03, xmax=0.32, ymin=0.0, ymax=0.38,
  xtick={0.05,0.10,0.15,0.20,0.25,0.30},
  xticklabels={5\%,10\%,15\%,20\%,25\%,30\%},
  ytick={0,0.05,0.10,0.15,0.20,0.25,0.30,0.35},
  legend pos=north west,
  legend columns=2,
]
\addplot[blue,thick,mark=*,solid] coordinates {
  (0.05,0.086)(0.10,0.141)(0.15,0.197)(0.20,0.225)(0.25,0.256)(0.30,0.286)
};
\addlegendentry{Proposed}
\addplot[red!80!black,thick,mark=square,dashed] coordinates {
  (0.05,0.079)(0.10,0.079)(0.15,0.080)(0.20,0.080)(0.25,0.072)(0.30,0.090)
};
\addlegendentry{Centralized}
\addplot[orange,thick,mark=triangle,dotted] coordinates {
  (0.05,0.149)(0.10,0.134)(0.15,0.145)(0.20,0.141)(0.25,0.136)(0.30,0.160)
};
\addlegendentry{Static DLT}
\addplot[green!60!black,thick,mark=diamond,dash dot] coordinates {
  (0.05,0.258)(0.10,0.220)(0.15,0.244)(0.20,0.245)(0.25,0.249)(0.30,0.243)
};
\addlegendentry{ZT Auth-Only}
\addplot[purple,thick,mark=+,loosely dashed] coordinates {
  (0.05,0.084)(0.10,0.138)(0.15,0.194)(0.20,0.222)(0.25,0.255)(0.30,0.285)
};
\addlegendentry{UAV+FL}
\end{axis}
\end{tikzpicture}
\caption{Attack resilience vs malicious fraction $f$ (5--30\%),
         at $\theta\!=\!0.40$. Proposed model increases monotonically;
         Centralized Trust plateaus at F1\,$\approx$\,0.08 due to
         partition-blocked revocation.}
\label{fig:attack}
\end{figure}

% ===========================================================================
\section{Segment-Optimal Threshold Calibration}
\label{sec:threshold}

\subsection{Steady-State Trust Analysis}

Setting $T_{ij}(t\!+\!\Delta t) = T_{ij}(t) = T^*_s$ in
Eq.~\eqref{eq:aging} yields the closed-form steady-state trust:
\begin{equation}
  T^*_s =
    \frac{(1-\alpha)\,\bar{s}}{1 - \alpha\,e^{-\lambda_s \Delta t_{\mathrm{eff}}}},
  \label{eq:steady}
\end{equation}
where $\Delta t_{\mathrm{eff}} = T_{\mathrm{round}} + \bar{d}_s/1000$.
Table~\ref{tab:lambda} lists $T^*_{\mathrm{legit}}$ (with
$\bar{s}\!=\!0.85$) and $T^*_{\mathrm{mal}}$ (with $\bar{s}\!=\!0.10$)
for each segment.

\begin{table}[t]
\centering
\caption{Steady-State Trust Values and Optimal Thresholds per Segment
         ($\alpha\!=\!0.70$, $T_{\mathrm{round}}\!=\!60$\,s)}
\label{tab:lambda}
\renewcommand{\arraystretch}{1.15}
\small
\begin{tabular}{lrrrr}
\toprule
Segment & $\lambda_s$ (s$^{-1}$) & $T^*_{\mathrm{legit}}$ &
          $T^*_{\mathrm{mal}}$ & $\theta^*_s$ \\
\midrule
GEO    & 0.005 & 0.530 & 0.062 & 0.296 \\
Ground & 0.015 & 0.356 & 0.042 & 0.199 \\
LEO    & 0.020 & 0.323 & 0.038 & 0.181 \\
UAV    & 0.050 & 0.264 & 0.031 & 0.148 \\
\bottomrule
\end{tabular}
\end{table}

\subsection{Round-Period Sensitivity}
\label{sec:tround}

Table~\ref{tab:tround_sensitivity} extends the steady-state analysis
across round durations for the Ground segment ($\lambda\!=\!0.015$\,s$^{-1}$,
ignoring propagation delay for clarity).
As $T_{\mathrm{round}}$ increases, the exponential factor
$e^{-\lambda_s \Delta t}$ decreases, reducing $T^*_s$ and tightening
the legitimate--malicious separation.
The global threshold property
$T^*_{\mathrm{mal}} < \theta^* < T^*_{\mathrm{legit}}$ is preserved
for all $T_{\mathrm{round}} > 0$; the analytically calibrated formula
Eq.~\eqref{eq:theta_opt} remains valid across all tested durations.

\begin{table}[t]
\centering
\caption{Analytical Sensitivity to Round Period
         (Ground: $\lambda\!=\!0.015$\,s$^{-1}$, $\alpha\!=\!0.70$)}
\label{tab:tround_sensitivity}
\renewcommand{\arraystretch}{1.1}
\small
\begin{tabular}{rrrr}
\toprule
$T_{\mathrm{round}}$ (s) & $T^*_{\mathrm{legit}}$ &
$T^*_{\mathrm{mal}}$ & $\theta^*_{\mathrm{Ground}}$ \\
\midrule
 30 & 0.461 & 0.054 & 0.257 \\
\textbf{60} & \textbf{0.357} & \textbf{0.042} & \textbf{0.199} \\
 90 & 0.312 & 0.037 & 0.174 \\
120 & 0.288 & 0.034 & 0.161 \\
\bottomrule
\end{tabular}
{\footnotesize Bold: simulated operating point. Values from
 Eq.~\eqref{eq:steady}; $T^*_{\mathrm{mal}}/T^*_{\mathrm{legit}}$
 ratio $\leq 0.12$ across all rows confirms robust separation.}
\end{table}

\subsection{Segment-Optimal Threshold Formula}

The threshold minimizing combined FPR and FNR for segment $s$:
\begin{equation}
  \theta^*_s =
    \frac{T^*_{s,\mathrm{legit}} + T^*_{s,\mathrm{mal}}}{2}.
  \label{eq:theta_opt}
\end{equation}
A global threshold $\theta^*\!=\!0.20$ simultaneously satisfies
$T^*_{s,\mathrm{mal}} < 0.20 < T^*_{s,\mathrm{legit}}$ for all
segments (minimum legitimate: 0.264~UAV; maximum malicious: 0.062~GEO).

\subsection{Empirical Confirmation (Exp.~7)}

Fig.~\ref{fig:threshold} and Table~\ref{tab:theta_sweep} present the
empirical sweep from $\theta\!=\!0.10$ to $0.50$ (30 runs each).
At $\theta^*\!=\!0.20$: \textbf{FPR\,=\,0.000, Precision\,=\,1.000,
F1\,=\,0.250}.
The FPR transition occurs sharply at $\theta\!\approx\!0.30$,
consistent with the analytical boundary at $T^*_{\mathrm{Ground}}$.
Recall is bounded by $\approx\!0.15$ because on-off attackers in
their compliant phase maintain trust above any fixed threshold by
design; detection latency (35.9~rounds) is the complementary metric.

\begin{figure}[t]
\centering
\begin{tikzpicture}
% (a) F1 and FPR vs theta
\begin{axis}[
  ieeehalf,
  at={(0,0)},
  xlabel={Threshold $\theta$},
  ylabel={F1},
  xmin=0.08, xmax=0.52, ymin=0.0, ymax=0.38,
  xtick={0.10,0.20,0.30,0.40,0.50},
  ytick={0,0.10,0.20,0.30},
  title={(a) F1 and FPR vs $\theta$},
  axis y line*=left,
  axis x line*=bottom,
  legend style={at={(0.02,0.98)},anchor=north west,font=\tiny},
]
\addplot[blue,thick,mark=o,solid,
         error bars/.cd,y dir=both,y explicit] coordinates {
  (0.10,0.199)+-(0,0.020)
  (0.15,0.244)+-(0,0.021)
  (0.20,0.250)+-(0,0.022)
  (0.25,0.254)+-(0,0.022)
  (0.30,0.300)+-(0,0.019)
  (0.35,0.284)+-(0,0.015)
  (0.40,0.236)+-(0,0.011)
  (0.45,0.237)+-(0,0.011)
  (0.50,0.236)+-(0,0.011)
};
\addlegendentry{F1}
\addplot[green!60!black,thin,dashed,mark=none] coordinates {
  (0.20,0)(0.20,0.38)};
\addlegendentry{$\theta^*\!=\!0.20$}
\end{axis}
\begin{axis}[
  ieeehalf,
  at={(0,0)},
  xmin=0.08, xmax=0.52, ymin=0.0, ymax=0.32,
  xtick={0.10,0.20,0.30,0.40,0.50},
  ytick={0,0.10,0.20,0.30},
  ylabel={FPR},
  ylabel style={red},
  yticklabel style={red},
  axis y line*=right,
  axis x line=none,
  legend style={at={(0.98,0.60)},anchor=east,font=\tiny},
]
\addplot[red,thick,mark=square,dashed] coordinates {
  (0.10,0.000)(0.15,0.000)(0.20,0.000)(0.25,0.000)
  (0.30,0.002)(0.35,0.108)(0.40,0.265)(0.45,0.275)(0.50,0.276)
};
\addlegendentry{FPR}
\end{axis}
% (b) Precision and Recall vs theta
\begin{axis}[
  ieeehalf,
  at={(0.53\columnwidth,0)},
  xlabel={Threshold $\theta$},
  ylabel={Score},
  xmin=0.08, xmax=0.52, ymin=0.0, ymax=1.10,
  xtick={0.10,0.20,0.30,0.40,0.50},
  ytick={0,0.20,0.40,0.60,0.80,1.00},
  title={(b) Precision and Recall vs $\theta$},
  legend pos=south east,
]
\addplot[green!60!black,thick,mark=triangle,solid] coordinates {
  (0.10,1.000)(0.15,1.000)(0.20,1.000)(0.25,1.000)
  (0.30,0.959)(0.35,0.354)(0.40,0.206)(0.45,0.204)(0.50,0.203)
};
\addlegendentry{Precision}
\addplot[orange!90!black,thick,mark=diamond,dashed] coordinates {
  (0.10,0.112)(0.15,0.140)(0.20,0.144)(0.25,0.147)
  (0.30,0.179)(0.35,0.238)(0.40,0.277)(0.45,0.283)(0.50,0.283)
};
\addlegendentry{Recall}
\addplot[green!60!black,thin,dashed,mark=none] coordinates {
  (0.20,0)(0.20,1.10)};
\end{axis}
\end{tikzpicture}
\caption{Threshold sensitivity (Exp.~7, 30 runs, $\pm$95\,\%\,CI).
         (a)~F1 (blue, left) and FPR (red, right) vs $\theta$.
         (b)~Precision and Recall vs $\theta$.
         Dotted green: analytically derived $\theta^*\!=\!0.20$
         (FPR\,=\,0, Precision\,=\,1.000).}
\label{fig:threshold}
\end{figure}

\begin{table}[t]
\centering
\caption{Threshold Sensitivity: Proposed Model at $f\!=\!20\%$
         (30 runs, mean values)}
\label{tab:theta_sweep}
\renewcommand{\arraystretch}{1.15}
\small
\begin{tabular}{rrrrr}
\toprule
$\theta$ & F1 & FPR & Precision & Recall \\
\midrule
0.10 & 0.199 & 0.000 & 1.000 & 0.112 \\
0.15 & 0.244 & 0.000 & 1.000 & 0.140 \\
\textbf{0.20} & \textbf{0.250} & \textbf{0.000} & \textbf{1.000} & \textbf{0.144} \\
0.25 & 0.254 & 0.000 & 1.000 & 0.147 \\
0.30 & 0.300 & 0.002 & 0.959 & 0.179 \\
0.35 & 0.284 & 0.108 & 0.354 & 0.238 \\
0.40 & 0.236 & 0.265 & 0.206 & 0.277 \\
0.45 & 0.237 & 0.275 & 0.204 & 0.283 \\
0.50 & 0.236 & 0.276 & 0.203 & 0.283 \\
\bottomrule
\end{tabular}
{\footnotesize Bold: analytically derived $\theta^*\!=\!0.20$.}
\end{table}

% ===========================================================================
\section{Discussion}
\label{sec:discussion}

\textbf{Detection latency is the decisive advantage.}
The $4.3\times$ speedup over Static DLT is mechanistically driven by
per-segment decay: $L_s\!=\!0.284$ (proposed) vs $L\!=\!0.70$
(Static DLT), meaning the proposed model's trust responses are
$7\times$ faster in the exponential sense.
This matters most for GEO-scale latency where stale trust from a
compromised satellite could persist for hours without decay.

\textbf{Deploy at $\theta^*\!=\!0.20$.}
Equation~\eqref{eq:theta_opt} provides a closed-form deployable formula.
Using $\theta\!=\!0.40$ (conservative experiment value) is demonstrably
sub-optimal for non-GEO segments, crossing above
$T^*_{\mathrm{Ground}}\!=\!0.356$ and causing FPR\,=\,26.8\,\%.
At $\theta\!=\!0.20$, zero false alarms and Precision\,=\,1.000.

\textbf{2048-bit keys are feasible.}
At $n\!=\!50$ contributors, overhead is 5{,}575\,ms (9.3\,\% of budget).
For denser clusters ($n\!>\!100$), batching into sub-clusters or
migrating to CKKS lattice-based HE reduces overhead further.

\textbf{Limitations.}
The simulation uses a random graph model rather than a 3GPP NR-NTN
channel model.
Consensus throughput is simulated without a networked validator
implementation.
The DP analysis uses standard $(\varepsilon,\delta)$ composition;
R\'{e}nyi DP~\cite{mironov2017} would yield tighter multi-round bounds.
Recall is bounded at $\approx\!0.15$ due to on-off adversarial evasion.

% ===========================================================================
\section{Conclusion}
\label{sec:conclusion}

This paper presented the first quantitative simulation-based evaluation
of a trustless-by-design security framework for 6G Non-Terrestrial
Networks.
Across eight experiments on a 1{,}000-node heterogeneous NTN
(30 Monte-Carlo runs each), the proposed framework achieved:
(i)~$4.3\times$ and $6.1\times$ detection latency improvement over
Static DLT and Centralized Trust;
(ii)~a $+0.130$ partition recovery F1 gain;
(iii)~near-linear $\mathcal{O}(N)$ scalability to 5{,}000 nodes;
(iv)~zero false positives at the analytically derived
$\theta^*\!=\!0.20$; and
(v)~NIST-compliant 2048-bit Paillier HE within the 60-second round
budget for $n\!\leq\!100$ contributors.

Future work should address hardware-in-the-loop validation with an
SDR-based LEO emulator, CKKS lattice HE for dense clusters, federated
ZTA policy adaptation, R\'{e}nyi DP multi-round accounting, and
repeated partition evaluation under GEO--LEO handover cycles.

% ===========================================================================
\appendix[Proof of Steady-State Trust Convergence]
\label{app:convergence}

The trust update map $f(x) = L_s x + (1\!-\!\alpha)\bar{s}$
on $[0,1]$, where $L_s = \alpha e^{-\lambda_s \Delta t}$,
is a contraction with Lipschitz constant $L_s < 1$.
By the Banach fixed-point theorem it converges from any initial value
to the unique fixed point Eq.~\eqref{eq:steady}.
Convergence to within 1\,\% of $T^*_s$ requires
$k \geq \lceil\ln(0.01)/\ln(L_s)\rceil$ rounds:
$k\!=\!8$ for GEO ($L_s\!=\!0.518$) and $k\!=\!2$ for UAV
($L_s\!=\!0.035$).
\QED

% ===========================================================================
\begin{thebibliography}{15}

\bibitem{giordani2020}
M.~Giordani, M.~Polese, M.~Mezzavilla, S.~Rangan, and M.~Zorzi,
``Toward 6G networks: Use cases and technologies,''
\textit{IEEE Commun. Mag.}, vol.~58, no.~3, pp.~55--61, Mar.~2020.

\bibitem{yang2024}
H.~Yang, Y.~Yu, Y.~Zhu, and J.~Li,
``Toward trustworthy 6G networks: A trust-based consensus scheme,''
\textit{IEEE Network}, vol.~38, no.~1, pp.~98--104, Jan.~2024.

\bibitem{nxgent}
C.~N\'{u}\~{n}ez-G\'{o}mez, V.~Garcia-Font, H.~Rif\`{a}-Pous,
M.~Asad, and J.~Solera,
``NxGenT: A decentralized trust system for B5G and 6G networks,''
\textit{Blockchain: Res. Appl.}, vol.~5, 2024.

\bibitem{kalla2022}
A.~Kalla, C.~de~Alwis, P.~Porambage, G.~G\"{u}r, and M.~Liyanage,
``A survey on the use of blockchain for future 6G networks,''
\textit{J. Ind. Inf. Integration}, vol.~27, p.~100317, 2022.

\bibitem{chen2024}
X.~Chen, W.~Feng, N.~Ge, and Y.~Zhang,
``Zero trust architecture for 6G security,''
\textit{IEEE Network}, vol.~38, no.~4, pp.~120--127, 2024.

\bibitem{ztisl}
S.~R.~Pokhrel,
``Orbital ZTA! Secure satellite communication networks with zero trust
architecture,''
in \textit{Proc. ACM SIGCOMM Workshop Security Space Syst.}, 2024,
pp.~15--22.

\bibitem{uav_fl}
A.~Ali and I.~K.~Shahbaz,
``SkyTrust: Blockchain-enhanced UAV security for NTNs with dynamic
trust and energy-aware consensus,''
\textit{arXiv preprint arXiv:2508.18735}, Aug.~2025.

\bibitem{tong2024}
Z.~Tong, J.~Wang, X.~Hou, J.~Chen, Z.~Jiao, and J.~Liu,
``Blockchain-based trustworthy and efficient hierarchical federated
learning for UAV-enabled IoT networks,''
\textit{IEEE Internet Things J.}, vol.~11, no.~21,
pp.~34270--34282, Nov.~2024.

\bibitem{aot}
Y.~Xiao, Q.~Du, W.~Cheng, P.~D.~Diamantoulakis, and
G.~K.~Karagiannidis,
``Age of trust (AoT): A continuous verification framework for wireless
networks,''
\textit{IEEE Commun. Lett.}, 2025, doi:~10.1109/LCOMM.2025.3549401.

\bibitem{khan2023}
M.~Khan, S.~A.~Hassan, D.~O.~Akande, S.~Muyeen, and A.~Mahmood,
``Blockchain-based architectures for 5G and 6G networks: A survey,''
\textit{IEEE Access}, vol.~11, pp.~20946--20973, 2023.

\bibitem{paillier1999}
P.~Paillier,
``Public-key cryptosystems based on composite degree residuosity
classes,''
in \textit{Proc. EUROCRYPT}, 1999, pp.~223--238.

\bibitem{dwork2014}
C.~Dwork and A.~Roth,
``The algorithmic foundations of differential privacy,''
\textit{Found. Trends Theor. Comput. Sci.}, vol.~9, nos.~3--4,
pp.~211--407, 2014.

\bibitem{mironov2017}
I.~Mironov,
``R\'{e}nyi differential privacy of the Gaussian mechanism,''
in \textit{Proc. IEEE CSF}, 2017, pp.~263--275.

\bibitem{kaul2012}
S.~Kaul, R.~Yates, and M.~Gruteser,
``Real-time status: How often should one update?''
in \textit{Proc. IEEE INFOCOM}, 2012, pp.~2731--2735.

\bibitem{zhou2019}
W.~Zhou, Y.~Jia, A.~Peng, Y.~Zhang, and P.~Liu,
``The effect of IoT new features on security and privacy,''
\textit{IEEE Internet Things J.}, vol.~6, no.~2, pp.~1606--1616, 2019.

\bibitem{conf_version}
C.~Palli, S.~Sarkar, and A.~Adhya,
``Trustless-by-design security framework for distributed IoT of 6G
non-terrestrial networks,''
in \textit{Proc. IEEE Veh. Technol. Conf. (VTC)}, 2026,
Paper no.~2026004467.

\end{thebibliography}

\end{document}
